\pagestyle{otFront}
\otPageTitre{Apprendre avec OptiTAB}{\OptionLivre · Terminale générale}
Ce manuel accompagne l'enseignement optionnel français de \OptionLivre.
Il associe un cours progressif, des méthodes expliquées, des exercices de trois niveaux et des solutions rédigées.
\begin{savoirfaire}[Un parcours en quatre temps]
\textbf{Comprendre.} Lire le cours et refaire l'exemple résolu, calculs masqués.\par
\textbf{S'entraîner.} Commencer par les exercices accessibles, puis travailler les applications.\par
\textbf{Chercher.} Réserver un vrai temps aux problèmes, aux conjectures et à leur justification.\par
\textbf{Consolider.} Comparer les raisonnements, noter l'erreur puis refaire la question deux jours plus tard.
\end{savoirfaire}
Les points colorés indiquent une difficulté croissante. Les renvois d'un exercice vers sa solution, et de la solution vers l'énoncé, sont cliquables.
Les six compétences se travaillent ensemble : chercher, modéliser, représenter, raisonner, calculer et communiquer.
\souscours{Lire les énoncés avec précision}
Une hypothèse est une information disponible ; une conclusion est ce qu'il faut établir.
Une représentation ou une simulation aide à conjecturer. Une preuve explique pourquoi l'énoncé reste vrai dans tous les cas annoncés.
Chaque problème exige une conclusion dans son contexte et, si nécessaire, une unité, une précision ou une limite du modèle.
\souscours{Données et outils}
Sauf mention contraire, les données des exercices sont fictives et choisies pour l'apprentissage. Les situations médicales, économiques et cryptographiques sont des modèles pédagogiques.
Une valeur exacte précède autant que possible son arrondi. La calculatrice et Python servent à explorer, à calculer et à contrôler ; les arguments mathématiques restent explicités.
\vfill
{\small OptiTAB · Édition numérique, octobre 2026. Textes, exercices et figures de cette édition : création éditoriale OptiTAB. Les programmes officiels restent la propriété de leurs auteurs institutionnels. Ce manuel indépendant n'est pas une publication du ministère et ne revendique aucune homologation.}
\clearpage
\otPageTitre{Sommaire}{Un itinéraire complet, des outils aux synthèses}
\noindent\makebox[9mm][l]{\otSemi\color{otNavy}01}\hyperref[chap-1]{Calculer dans les complexes\otPoints\pageref*{chap-1}}\par\vspace{2.6mm}
\noindent\makebox[9mm][l]{\otSemi\color{otNavy}02}\hyperref[chap-2]{Module, arguments et plan complexe\otPoints\pageref*{chap-2}}\par\vspace{2.6mm}
\noindent\makebox[9mm][l]{\otSemi\color{otNavy}03}\hyperref[chap-3]{Exponentielle imaginaire et trigonométrie\otPoints\pageref*{chap-3}}\par\vspace{2.6mm}
\noindent\makebox[9mm][l]{\otSemi\color{otNavy}04}\hyperref[chap-4]{Équations polynomiales\otPoints\pageref*{chap-4}}\par\vspace{2.6mm}
\noindent\makebox[9mm][l]{\otSemi\color{otNavy}05}\hyperref[chap-5]{Géométrie et racines de l’unité\otPoints\pageref*{chap-5}}\par\vspace{2.6mm}
\noindent\makebox[9mm][l]{\otSemi\color{otNavy}06}\hyperref[chap-6]{Divisibilité et congruences\otPoints\pageref*{chap-6}}\par\vspace{2.6mm}
\noindent\makebox[9mm][l]{\otSemi\color{otNavy}07}\hyperref[chap-7]{PGCD, Bézout et Gauss\otPoints\pageref*{chap-7}}\par\vspace{2.6mm}
\noindent\makebox[9mm][l]{\otSemi\color{otNavy}08}\hyperref[chap-8]{Nombres premiers et petit théorème de Fermat\otPoints\pageref*{chap-8}}\par\vspace{2.6mm}
\noindent\makebox[9mm][l]{\otSemi\color{otNavy}09}\hyperref[chap-9]{Matrices et systèmes linéaires\otPoints\pageref*{chap-9}}\par\vspace{2.6mm}
\noindent\makebox[9mm][l]{\otSemi\color{otNavy}10}\hyperref[chap-10]{Puissances et suites matricielles\otPoints\pageref*{chap-10}}\par\vspace{2.6mm}
\noindent\makebox[9mm][l]{\otSemi\color{otNavy}11}\hyperref[chap-11]{Graphes et matrices d’adjacence\otPoints\pageref*{chap-11}}\par\vspace{2.6mm}
\noindent\makebox[9mm][l]{\otSemi\color{otNavy}12}\hyperref[chap-12]{Chaînes de Markov à deux ou trois états\otPoints\pageref*{chap-12}}\par\vspace{2.6mm}
\vspace{4mm}\repere{dossiers}{Dossiers et évaluations de synthèse}
\repere{corriges-debut}{Corrigés détaillés}
\repere{concordance}{Concordance avec le programme et sources}
\repere{index-methodes}{Index des méthodes et formules}
\clearpage
\otPageTitre{Faire le point}{Diagnostic de départ · À reprendre après quelques semaines}

\begin{questions}
\item Développer $(x+2)^3$ et factoriser $x^3-1$.
\item Donner $\cos(\pi/3)$, $\sin(\pi/3)$ et $\cos(\pi/2)$.
\item Résoudre $2x+3y=7$, $x-y=1$.
\item Écrire la division euclidienne de $47$ par $6$ ; donner le PGCD de $18$ et $30$.
\item Une probabilité de transition de $A$ vers $B$ vaut $0{,}3$. Quel est le poids de $A$ vers $A$ s'il n'y a que deux états ?
\item Pour $u_0=1$, $u_{n+1}=2u_n+1$, donner les quatre premiers termes.
\end{questions}
\begin{savoirfaire}[Réponses et orientation]
1. $x^3+6x^2+12x+8$ ; $(x-1)(x^2+x+1)$.
2. $1/2$, $\sqrt3/2$, $0$.
3. $x=2,y=1$.
4. $47=6\times7+5$, PGCD $6$.
5. $0{,}7$ : les transitions sortant de $A$ totalisent $1$.
6. $1,3,7,15$.
Les chapitres 1 à 5 prolongent l'algèbre et la trigonométrie ; les chapitres 6 à 8 structurent l'arithmétique ; les chapitres 9 à 12 développent matrices, graphes et chaînes de Markov.
\end{savoirfaire}
Les mathématiques expertes s'articulent avec la spécialité de Terminale. Les rappels sont fournis lorsqu'ils servent les nouveaux objets, mais le manuel de spécialité reste utile pour les suites, l'analyse et les probabilités.

\clearpage\otPageTitre{Mon parcours}{Cocher après une réussite autonome, puis après révision}
\begin{tabularx}{\linewidth}{@{}p{8mm}Yccc@{}}\toprule
& Chapitre & Cours & Exercices & Révision\\\midrule
1 & Calculer dans les complexes & $\square$ & $\square$ & $\square$ \\[2mm]
2 & Module, arguments et plan complexe & $\square$ & $\square$ & $\square$ \\[2mm]
3 & Exponentielle imaginaire et trigonométrie & $\square$ & $\square$ & $\square$ \\[2mm]
4 & Équations polynomiales & $\square$ & $\square$ & $\square$ \\[2mm]
5 & Géométrie et racines de l’unité & $\square$ & $\square$ & $\square$ \\[2mm]
6 & Divisibilité et congruences & $\square$ & $\square$ & $\square$ \\[2mm]
7 & PGCD, Bézout et Gauss & $\square$ & $\square$ & $\square$ \\[2mm]
8 & Nombres premiers et petit théorème de Fermat & $\square$ & $\square$ & $\square$ \\[2mm]
9 & Matrices et systèmes linéaires & $\square$ & $\square$ & $\square$ \\[2mm]
10 & Puissances et suites matricielles & $\square$ & $\square$ & $\square$ \\[2mm]
11 & Graphes et matrices d’adjacence & $\square$ & $\square$ & $\square$ \\[2mm]
12 & Chaînes de Markov à deux ou trois états & $\square$ & $\square$ & $\square$ \\[2mm]
\bottomrule\end{tabularx}\par\vspace{5mm}Une case indique une réussite sans aide. Si un exercice résiste, revenir au cours et au premier exemple avant de reprendre le problème.\clearpage\pagestyle{otMain}